
\section{DCN line extraction, background subtraction and hyperfine fitting}
\label{appendix:dcn_fitting}

 
We use \texttt{spectral-cube}\footnote{\url{https://spectral-cube.readthedocs.io/en/latest/}} \citep{Ginsburg2019b} to extract the DCN (3-2) emission from the Band 6 spectral window 1 cubes. First, we shift the cubes into the  frequency rest frame of the DCN (3-2) transition, and then for each protocluster we extract a subcube with a 50~\kms width centered on the \vlsr of the protocluster. We then extract a core-averaged, background subtracted DCN (3-2) spectrum for each continuum core. To do this, we use the smoothed core catalog from \citet{louvet22} to obtain the positions and core sizes for the 595 continuum cores over the 15 ALMA-IMF protoclusters. We extract the core-averaged DCN (3-2) spectrum from all pixels within an elliptical aperture with major (minor) axis length twice that of the continuum source major (minor) FWHM. We subtract from this an average background spectrum taken from within an elliptical annulus with inner major (minor) axis size equal to the size of the spectral aperture (i.e., twice the continuum source major (minor) FWHM) and outer radius size 1.5x larger (i.e., three times the continuum source major (minor) FWHM)  %Following the same framework as the continuum core extraction, we also subtract the background emission from these ON spectra. 
As nearby cores can overlap with the background annulus, we exclude those pixels from neighboring cores that spatially overlap with the background annulus, creating a background ``crown''. An example of the core and background annulus/crown definition can be seen in \cref{appendix:dcn_zoom_masks}.  
The resulting core-averaged, background-subtracted spectrum for each core is then fitted using \texttt{PySpecKit} \citep{Ginsburg2011} with an adapted single component hyperfine model (HFS), to account for the hyperfine structure in the DCN (3-2) line. The values for the specified hyperfine model are taken from the CDMS catalog to define the frequencies and relative weights of the hyperfine lines \citep{Endres2016,muller2001,muller2005} and are presented in \cref{tabappendix:hyperfine}. An example of the positions of the hyperfine components overlaid on an example DCN spectrum is presented in \cref{appendix:dcn_hyperfineexample}.
\label{sec:hyperfinemodel}

 
\subsection{Single- and Complex-type spectral classification}

To classify a spectrum as either detection or non-detection, we require a SNR of at least 4~$\sigma$ between its peak and its estimated noise, and a velocity dispersion of $> 0.2$~\kms. The noise in the spectrum is estimated from the MAD\footref{note:mad} in 30 continuous channels ($\sim$10\,\kms) from either the lower or upper part of the spectrum (shaded pink regions in the DCN spectra), that were identified by eye to be in a line free part of the spectrum. We then split the detected spectra in Single- or Complex-type. This is done manually through a visual inspection of the DCN (3-2) spectra after fitting. If multiple components are present in the spectrum, we also  classify them as Complex-type and  remove them from the analysis. The spectra for each protocluster are shown in \cref{appendix:spectral}. 

\subsection{Comparison of the background and non-background subtracted DCN spectra}

In Figures \ref{figspectra:dcnspectra_split_G338_s}, and \ref{appendicfig:g388_on}, we show the DCN (3-2) detected Single- and Complex-type spectra identified when performing background-subtraction, and non-background subtraction, respectively, for G338.93. In Figure \ref{g338:off}, for reference we also show the background spectra. We find that some of the "absorption features" in the background subtracted spectra are a result of the background subtraction (e.g., core 13, and 29). The background-subtraction also results in some cores with initially Complex-type spectra due to the background contamination becoming Single-type spectra (e.g., core 16). If we consider the detection statistics for the background and non background subtraction they are similar in G338.93, 21 and 18, respectively, for Single-type spectra (with some small differences between which cores are Single-type in the two sets), while for the Complex-type spectra both have 8. Furthermore, the average \vlsr extracted from all cores is the same (-61~\kms) when considering either background- or non background-subtraction. The average linewidth of the fits, however, is smaller ($\sim$1.6~\kms) for the background-subtracted compared with ($\sim$2.1~\kms) for the non-background subtracted, which is due to the removal of the background contribution. These background-subtracted spectra are likely to be more representative of the gas associated with the cores. In \cref{tabappendix:dcnfits_average}, we show the same table as in \cref{tab:dcnfits_average} but from fits to the non-background subtracted spectra for all protoclusters. This highlights that the general trends across the sample are the same regardless of whether the background-subtracted or non-background subtracted spectra are used.


\begin{table*}[htbp!]
\centering
%\small
%    \setlength{\tabcolsep}{2pt}
\begin{threeparttable}[c]
\caption{Characteristic parameters of the DCN (3-2) hyperfine fits for the non-background subtracted}
%(using global measurements)}
\label{tabappendix:dcnfits_average}
 \begin{tabular}{lllllllllllll}
\hline \noalign {\smallskip}
%Protocluster  & \multicolumn{3}{c}{spw0 }  & \multicolumn{3}{c}{spw1}  & \multicolumn{3}{c}{spw2 }
%    & \multicolumn{3}{c}{spw3} & \\ %& \lbol & \tdust 
%cloud name\tnote{1}    & min & max & rest   & min & max & rest  & min & max & rest  & min & max & rest   \\
Protocluster         & Number\tnote{1}  & Number\tnote{2}   & Detection\tnote{2} & Velocity\tnote{3} & \vlsr\tnote{3}   & \vlsr\tnote{3}    & Linewidth\tnote{3}  & Linewidth\tnote{3}   &  Linewidth\tnote{3}  \\ 
cloud name & of Cores & detected    & rate [\%]   & range    & mean   & std     & mean  & median &  std \\
&  &(complex)        &  (complex)      & [km~s$^{-1}$]   & [km~s$^{-1}$] & [km~s$^{-1}$]  & [km~s$^{-1}$]& [km~s$^{-1}$] & [km~s$^{-1}$] \\

\hline \noalign {\smallskip}
\multicolumn{10}{c}{Young} \\
G327.29 & 32 & 18 (23)&  56 (72)\% & 5.3 & -44.9 & 1.9  & 1.7 & 1.7 & 0.5 \\
G328.25 & 11 & 3 (4)&  27 (36)\% & 1.4 & -43.1 & 0.6  & 1.9 & 1.6 & 0.9 \\
G338.93 & 42 & 21 (29)&  50 (69)\% & 6.9 & -60.9 & 2.0  & 2.1 & 2.0 & 0.7 \\
G337.92 & 22 & 10 (17)&  45 (77)\% & 4.0 & -39.2 & 1.3  & 2.1 & 2.0 & 1.1 \\
W43-MM1 & 70 & 32 (51)&  46 (73)\% & 7.5 & 97.2 & 2.1  & 1.9 & 1.6 & 1.1 \\
W43-MM2 & 40 & 23 (30)&  58 (75)\% & 4.3 & 90.8 & 1.3  & 1.9 & 1.6 & 0.9 \\
Average &  36 & 18(26) & 49 (71)\% &  4.9 &  & 1.8  & 1.9 & 1.7 & 0.9 \\
Total &  217 & 107 (154)  \\
\hline
\multicolumn{10}{c}{Intermediate} \\
G351.77 & 18 & 5 (14)&  28 (78)\% & 5.9 & -4.5 & 2.1  & 2.0 & 1.6 & 0.6 \\
G353.41 & 45 & 17 (22)&  38 (49)\% & 7.4 & -17.3 & 1.8  & 1.7 & 1.5 & 0.8 \\
G008.67 & 19 & 8 (13)&  42 (68)\% & 6.6 & 35.8 & 2.0  & 2.0 & 1.9 & 0.6 \\
W43-MM3 & 36 & 12 (18)&  33 (50)\% & 2.6 & 93.1 & 0.8  & 1.8 & 2.0 & 0.7 \\
W51-E & 23 & 12 (21)&  52 (91)\% & 11.1 & 56.8 & 3.9  & 2.0 & 1.6 & 1.0 \\
Average &  28 & 11(18) & 38 (62)\% &  6.7 &  & 2.8  & 1.8 & 1.6 & 0.8 \\
Total &  141 & 54 (88)  \\
\hline
\multicolumn{10}{c}{Evolved} \\
G333.60 &  52 & 36 (42) &  69 (81)\% & 10.0 & -47.8 & 1.9  & 1.7 & 1.7 & 0.5 \\
G010.62 &  42 & 28 (37) &  67 (88)\% & 9.3 & -2.5 & 2.1  & 1.4 & 1.3 & 0.6 \\
G012.80 &  46 & 38 (40) &  83 (87)\% & 6.4 & 36.1 & 1.4  & 1.4 & 1.2 & 0.6 \\
W51-IRS2 &  97 & 69 (84) &  71 (87)\% & 18.3 & 61.9 & 2.9  & 1.5 & 1.3 & 0.8 \\
Average &  59 & 43 (51) & 72 (86)\% &  11.0 &  & 2.5  & 1.5 & 1.3 & 0.7 \\
Total &  237 & 171 (203)  \\
\hline \noalign {\smallskip}
\end{tabular}
\begin{tablenotes}
\item[1] The number of continuum cores are taken from the smoothed core catalog from \citet{louvet22}. %We note that 2 cores overlap in W51-E and W51-IRS2 regions and 2 cores overlap in W43-MM2/3 regions (we note these in \cref{appendix:dcn_fitting}), however, as they are either non-detections or Complex-type spectra we do not explicitly assign them to a given field. 
\item[2] The total count/percentages in the brackets include both Single- and Complex-type DCN (3-2) detections. Where spectra that are described as Complex-type spectra and are not well fit by a single component. We provide the percentages and number of detections, but cores with a Complex-type spectrum are not included in the estimates of the Velocity range, \vlsr, or linewidth.
\item[3] Only Single-type DCN spectra (fit with a single component) and listed as Single in the Tables of the DCN fits (e.g. \cref{tab:coretables_g338}) are included in the estimates of the Velocity ranges, \vlsr and linewidth for each protocluster. 
%
%
\end{tablenotes}
\end{threeparttable}
\end{table*}


\begin{figure*}
    \centering
    \includegraphics[width=0.44\textwidth]{Appendix/appendix_figures/appendix_coresoverlaid_G333_core14.png}
    \includegraphics[width=0.44\textwidth]{Appendix/appendix_figures/appendix_coresoverlaid_G333_core45.png}
    \includegraphics[width=0.48\textwidth]{Appendix/appendix_figures/appendix_core_14_mask_G333.png}
    \includegraphics[width=0.48\textwidth]{Appendix/appendix_figures/appendix_core_45_mask_G333.png}
    \caption{ Zoom in toward cores 14 (top-left panel), and 45 (top-right panel) overlaid on the DCN moment 0 map in the protocluster G333.60.  The red ellipses are 2 times the major (minor) FWHM of the continuum cores. In the bottom panels we show examples of the masks (shaded blue areas), and their respective annulus (for core 14) or crown (for core 45), that is used to estimate the core-averaged, background-subtracted DCN (3-2) spectra. The background annulus are defined as 2-3 times the core major (minor) FWHMs from \citet{louvet22}, with pixels from overlapping cores being masked out (i.e., the resulting crown in core 45). \label{appendix:dcn_zoom_masks}} 
\end{figure*}

\begin{figure}
    \centering
    \includegraphics[width=0.44\textwidth]{Appendix/appendix_figures/appendix-sectra_14_g333_zoom.png}
    \caption{ Example of the DCN (3-2) core-averaged, background-subtracted spectrum extracted for core 14, in G333.60 with the DCN hyperfine model (see \cref{tabappendix:hyperfine} overlaid in red vertical dashed lines and are set relative to the central core \vlsr. \label{appendix:dcn_hyperfineexample}}
\end{figure}

\begin{figure*}

    \centering
\includegraphics[width=0.24\textwidth]{Appendix/appendix_figures/spectra_plots/G338_ON_spectra/G338.93_core_1_hyperfine_fits_on.png}
\includegraphics[width=0.24\textwidth]{Appendix/appendix_figures/spectra_plots/G338_ON_spectra/G338.93_core_2_hyperfine_fits_on.png}
\includegraphics[width=0.24\textwidth]{Appendix/appendix_figures/spectra_plots/G338_ON_spectra/G338.93_core_5_hyperfine_fits_on.png}
\includegraphics[width=0.24\textwidth]{Appendix/appendix_figures/spectra_plots/G338_ON_spectra/G338.93_core_6_hyperfine_fits_on.png}
\includegraphics[width=0.24\textwidth]{Appendix/appendix_figures/spectra_plots/G338_ON_spectra/G338.93_core_8_hyperfine_fits_on.png}
\includegraphics[width=0.24\textwidth]{Appendix/appendix_figures/spectra_plots/G338_ON_spectra/G338.93_core_9_hyperfine_fits_on.png}
\includegraphics[width=0.24\textwidth]{Appendix/appendix_figures/spectra_plots/G338_ON_spectra/G338.93_core_10_hyperfine_fits_on.png}
\includegraphics[width=0.24\textwidth]{Appendix/appendix_figures/spectra_plots/G338_ON_spectra/G338.93_core_11_hyperfine_fits_on.png}
\includegraphics[width=0.24\textwidth]{Appendix/appendix_figures/spectra_plots/G338_ON_spectra/G338.93_core_13_hyperfine_fits_on.png}
\includegraphics[width=0.24\textwidth]{Appendix/appendix_figures/spectra_plots/G338_ON_spectra/G338.93_core_17_hyperfine_fits_on.png}
\includegraphics[width=0.24\textwidth]{Appendix/appendix_figures/spectra_plots/G338_ON_spectra/G338.93_core_19_hyperfine_fits_on.png}
\includegraphics[width=0.24\textwidth]{Appendix/appendix_figures/spectra_plots/G338_ON_spectra/G338.93_core_25_hyperfine_fits_on.png}
\includegraphics[width=0.24\textwidth]{Appendix/appendix_figures/spectra_plots/G338_ON_spectra/G338.93_core_28_hyperfine_fits_on.png}
\includegraphics[width=0.24\textwidth]{Appendix/appendix_figures/spectra_plots/G338_ON_spectra/G338.93_core_29_hyperfine_fits_on.png}
\includegraphics[width=0.24\textwidth]{Appendix/appendix_figures/spectra_plots/G338_ON_spectra/G338.93_core_31_hyperfine_fits_on.png}
\includegraphics[width=0.24\textwidth]{Appendix/appendix_figures/spectra_plots/G338_ON_spectra/G338.93_core_33_hyperfine_fits_on.png}
\includegraphics[width=0.24\textwidth]{Appendix/appendix_figures/spectra_plots/G338_ON_spectra/G338.93_core_41_hyperfine_fits_on.png}
%\includegraphics[width=0.24\textwidth]{Appendix/appendix_figures/spectra_plots/G338_ON_spectra/G338.93_core_42_hyperfine_fits_on.png}
\includegraphics[width=0.24\textwidth]{Appendix/appendix_figures/spectra_plots/G338_ON_spectra/G338.93_core_43_hyperfine_fits_on.png}
\includegraphics[width=0.24\textwidth]{Appendix/appendix_figures/spectra_plots/G338_ON_spectra/G338.93_core_45_hyperfine_fits_on.png}
\includegraphics[width=0.24\textwidth]{Appendix/appendix_figures/spectra_plots/G338_ON_spectra/G338.93_core_46_hyperfine_fits_on.png}
\includegraphics[width=0.24\textwidth]{Appendix/appendix_figures/spectra_plots/G338_ON_spectra/G338.93_core_47_hyperfine_fits_on.png}

\caption{The Single-type non-background subtracted spectra in the young protocluster G338.93. They are classified in the same way as the background-subtracted spectra. The associated continuum core number is given in the top left of each panel (the core numbering is taken from \citet{louvet22}) the core \vlsr (V) and velocity dispersion ($\sigma$) both in units of \kms from the HSF fit are given in the top right. The pink shaded region represents the part of the spectrum used to estimate the MAD noise. \label{appendicfig:g388_on}}
\end{figure*}


\begin{figure*}\ContinuedFloat
    \centering
\includegraphics[width=0.24\textwidth]{Appendix/appendix_figures/spectra_plots/G338_ON_spectra/G338.93_core_3_hyperfine_fits_on.png}
\includegraphics[width=0.24\textwidth]{Appendix/appendix_figures/spectra_plots/G338_ON_spectra/G338.93_core_4_hyperfine_fits_on.png}
\includegraphics[width=0.24\textwidth]{Appendix/appendix_figures/spectra_plots/G338_ON_spectra/G338.93_core_14_hyperfine_fits_on.png}
\includegraphics[width=0.24\textwidth]{Appendix/appendix_figures/spectra_plots/G338_ON_spectra/G338.93_core_16_hyperfine_fits_on.png}
\includegraphics[width=0.24\textwidth]{Appendix/appendix_figures/spectra_plots/G338_ON_spectra/G338.93_core_22_hyperfine_fits_on.png}
\includegraphics[width=0.24\textwidth]{Appendix/appendix_figures/spectra_plots/G338_ON_spectra/G338.93_core_23_hyperfine_fits_on.png}
\includegraphics[width=0.24\textwidth]{Appendix/appendix_figures/spectra_plots/G338_ON_spectra/G338.93_core_27_hyperfine_fits_on.png}
\includegraphics[width=0.24\textwidth]{Appendix/appendix_figures/spectra_plots/G338_ON_spectra/G338.93_core_40_hyperfine_fits_on.png}
\caption{Continued: The Complex-type, non-background subtracted spectra towards the young protocluster G338.93. \label{appendicfig:g388_on_c}}

\end{figure*}

\begin{figure*}

    \centering
\includegraphics[width=0.24\textwidth]{Appendix/appendix_figures/spectra_plots/G338_OFF_spectra/G338.93_core_1_hyperfine_fits_off.png}
\includegraphics[width=0.24\textwidth]{Appendix/appendix_figures/spectra_plots/G338_OFF_spectra/G338.93_core_2_hyperfine_fits_off.png}
\includegraphics[width=0.24\textwidth]{Appendix/appendix_figures/spectra_plots/G338_OFF_spectra/G338.93_core_5_hyperfine_fits_off.png}
\includegraphics[width=0.24\textwidth]{Appendix/appendix_figures/spectra_plots/G338_OFF_spectra/G338.93_core_6_hyperfine_fits_off.png}
\includegraphics[width=0.24\textwidth]{Appendix/appendix_figures/spectra_plots/G338_OFF_spectra/G338.93_core_8_hyperfine_fits_off.png}
\includegraphics[width=0.24\textwidth]{Appendix/appendix_figures/spectra_plots/G338_OFF_spectra/G338.93_core_9_hyperfine_fits_off.png}
\includegraphics[width=0.24\textwidth]{Appendix/appendix_figures/spectra_plots/G338_OFF_spectra/G338.93_core_10_hyperfine_fits_off.png}
\includegraphics[width=0.24\textwidth]{Appendix/appendix_figures/spectra_plots/G338_OFF_spectra/G338.93_core_11_hyperfine_fits_off.png}
\includegraphics[width=0.24\textwidth]{Appendix/appendix_figures/spectra_plots/G338_OFF_spectra/G338.93_core_13_hyperfine_fits_off.png}
\includegraphics[width=0.24\textwidth]{Appendix/appendix_figures/spectra_plots/G338_OFF_spectra/G338.93_core_17_hyperfine_fits_off.png}
\includegraphics[width=0.24\textwidth]{Appendix/appendix_figures/spectra_plots/G338_OFF_spectra/G338.93_core_19_hyperfine_fits_off.png}
\includegraphics[width=0.24\textwidth]{Appendix/appendix_figures/spectra_plots/G338_OFF_spectra/G338.93_core_25_hyperfine_fits_off.png}
\includegraphics[width=0.24\textwidth]{Appendix/appendix_figures/spectra_plots/G338_OFF_spectra/G338.93_core_28_hyperfine_fits_off.png}
\includegraphics[width=0.24\textwidth]{Appendix/appendix_figures/spectra_plots/G338_OFF_spectra/G338.93_core_29_hyperfine_fits_off.png}
%\includegraphics[width=0.24\textwidth]{Appendix/appendix_figures/spectra_plots/G338_OFF_spectra/G338.93_core_30_hyperfine_fits_off.png}
\includegraphics[width=0.24\textwidth]{Appendix/appendix_figures/spectra_plots/G338_OFF_spectra/G338.93_core_31_hyperfine_fits_off.png}
\includegraphics[width=0.24\textwidth]{Appendix/appendix_figures/spectra_plots/G338_OFF_spectra/G338.93_core_33_hyperfine_fits_off.png}
\includegraphics[width=0.24\textwidth]{Appendix/appendix_figures/spectra_plots/G338_OFF_spectra/G338.93_core_41_hyperfine_fits_off.png}
\includegraphics[width=0.24\textwidth]{Appendix/appendix_figures/spectra_plots/G338_OFF_spectra/G338.93_core_43_hyperfine_fits_off.png}
\includegraphics[width=0.24\textwidth]{Appendix/appendix_figures/spectra_plots/G338_OFF_spectra/G338.93_core_45_hyperfine_fits_off.png}
\includegraphics[width=0.24\textwidth]{Appendix/appendix_figures/spectra_plots/G338_OFF_spectra/G338.93_core_46_hyperfine_fits_off.png}
\includegraphics[width=0.24\textwidth]{Appendix/appendix_figures/spectra_plots/G338_OFF_spectra/G338.93_core_47_hyperfine_fits_off.png}

\caption{The background spectra which are the average spectra from the background annulus towards the positions of the Single-type cores shown in \cref{appendicfig:g388_on} in the young protocluster G338.93. The associated core numbers are provided in the top left of each panel. \label{g338:off}}
\end{figure*}


\begin{figure*}\ContinuedFloat
    \centering


\includegraphics[width=0.24\textwidth]{Appendix/appendix_figures/spectra_plots/G338_OFF_spectra/G338.93_core_3_hyperfine_fits_off.png}
\includegraphics[width=0.24\textwidth]{Appendix/appendix_figures/spectra_plots/G338_OFF_spectra/G338.93_core_4_hyperfine_fits_off.png}
\includegraphics[width=0.24\textwidth]{Appendix/appendix_figures/spectra_plots/G338_OFF_spectra/G338.93_core_14_hyperfine_fits_off.png}
\includegraphics[width=0.24\textwidth]{Appendix/appendix_figures/spectra_plots/G338_OFF_spectra/G338.93_core_16_hyperfine_fits_off.png}
\includegraphics[width=0.24\textwidth]{Appendix/appendix_figures/spectra_plots/G338_OFF_spectra/G338.93_core_22_hyperfine_fits_off.png}
\includegraphics[width=0.24\textwidth]{Appendix/appendix_figures/spectra_plots/G338_OFF_spectra/G338.93_core_23_hyperfine_fits_off.png}
\includegraphics[width=0.24\textwidth]{Appendix/appendix_figures/spectra_plots/G338_OFF_spectra/G338.93_core_27_hyperfine_fits_off.png}
\includegraphics[width=0.24\textwidth]{Appendix/appendix_figures/spectra_plots/G338_OFF_spectra/G338.93_core_16_hyperfine_fits_off.png}
\includegraphics[width=0.24\textwidth]{Appendix/appendix_figures/spectra_plots/G338_OFF_spectra/G338.93_core_40_hyperfine_fits_off.png}
\caption{Continued: The background spectra which are the average spectra from the background annulus towards the positions of the Complex-type cores shown in \cref{appendicfig:g388_on_c} in the young protocluster G338.93.}

\end{figure*}

\begin{table}[htbp!]
\centering
\caption{DCN hyperfine model}
\label{tabappendix:hyperfine}
 \begin{tabular}{llll}
\hline \noalign {\smallskip}
Hyperfine    & Frequency   & Relative Velocity &  Relative   \\ 
transition             &  [GHz]      &   [kms$^{-1}$]&        Line Strength  \\
\hline \noalign {\smallskip}
f1  &    217.236999  & 2.12      & 7   \\
f2  &    217.238300  & 0.328     & 5 \\
f3  &    217.238538  & 0.0       & 21 \\
f4  &    217.238555  &  -0.02373 & 7 \\
f5  &    217.238612  & -0.10239  &  9 \\
f6  &    217.239079  & -0.7468   & 5  \\
f7  &    217.240622  &  -2.8762  & 5\\
\hline \noalign {\smallskip}
\end{tabular}
\end{table}